% chapters/17/rangkuman.tex
\bagian{Rangkuman}

\begin{rangkumanbox}
\begin{itemize}
  \item Kompetensi mata kuliah ini disusun sebagai satu rantai: dua CPL di
        puncaknya, empat CPMK di bawahnya, tiga belas Sub-CPMK di bawahnya
        lagi, dan tujuh belas bab di dasarnya. Setiap tingkat menjawab
        pertanyaan yang berbeda.
  \item \CPL{03} memperoleh 57 persen dan \CPL{04} memperoleh 43 persen dari
        bobot mata kuliah ini. \CPMK{2} memiliki bobot terbesar, yaitu 32
        persen.
  \item Nilai akhir disusun dari enam komponen, dengan Proyek Akhir sebagai
        bobot terbesar. Proyek Akhir diberi bobot terbesar bukan karena paling
        melelahkan, melainkan karena ia satu-satunya instrumen yang mengukur
        kemampuan menggabungkan bagian-bagian.
  \item Setiap Sub-CPMK diukur oleh sekurang-kurangnya dua instrumen, sebab
        satu instrumen dapat keliru mengukur.
  \item Ketercapaian CPMK dinilai pada tingkat kelas, dengan ambang 75 persen
        mahasiswa bernilai 70 atau lebih. Apabila ambang tidak tercapai, yang
        diperbaiki adalah pelaksanaan perkuliahan, bukan mahasiswanya.
  \item Kelulusan menuntut nilai akhir sekurang-kurangnya 60, kehadiran
        sekurang-kurangnya 75 persen, dan penyerahan seluruh kelengkapan
        Proyek Akhir.
  \item Menguasai setiap bab belum tentu berarti menguasai permainan secara
        utuh. Kemampuan menggabungkan bagian-bagian adalah kompetensi
        tersendiri, dan dinilai secara terpisah.
  \item Lapisan permainan saling bergantung dari bawah ke atas, sehingga
        penelusuran kesalahan selalu bergerak ke bawah, bukan ke atas.
  \item Refleksi yang bernilai memuat keadaan yang konkret, tindakan yang
        diambil, dan hasil yang dapat diperiksa. Refleksi yang hanya memuat
        kesan tidak dapat ditindaklanjuti, baik oleh penilai maupun oleh Anda
        sendiri.
  \item Kebiasaan yang berulang pada seluruh bab ini adalah mengukur, bukan
        menduga: waktu satu bingkai, jangkauan lompatan, kesulitan level, dan
        kemajuan belajar semuanya diukur dengan cara yang dapat diperiksa
        ulang.
\end{itemize}
\end{rangkumanbox}
